\newpage
\section{Описание алгоритмов реализации методики}
\label{part:alghoritm_intro}


Несмотря на обширность множества в России приморских городов, выбор конкретного из них для меня был прост.
Самый крупный порт России, город-герой и город труженик, зажатый окружением курортных поселков -- город Новороссийск.
Руководствуясь его полуофициальным слоганом: \enquote{Берег России начинается здесь!} начнем анализ прибрежных территорий именно с него!

Текущая версия исполнения предлагаемой методики в целом согласуется с заявленной функциональностью описанной в разделе~\ref{part:methods_intro}.
Однако, на текущем этапе проработки задача была решена в контексте выполнения единственной задачи -- оценки морского побережья по предложенным в разделе~\ref{sec:wer_calculation} индексам WER.

Исходя из этого, в разработанные алгоритмы не были включены опции добавления и обработки внешних данных (геологии, топографии, экономики и~т.\,п.) в связи с отсутствием необходимости в них при анализе волновой уязвимости берега.
Таким образом, в текущей архитектуре решения отсутствует необходимость в выделении отдельной сущности \enquote{Задания} (заявленного в резделе~\ref{sec:task_intro} на стр.~\pageref{sec:task_end}), так как оно пока выполняется безальтернативно и, соответственно, отсутствует фабричная архитектура его формирования.
В перспективе это ограничение будет исправлено в следующих версиях.

Определенное несовершенство реализованной программы также заключается в том, что пока все хранение информации выполняется в файловом виде (форматах \texttt{GeoJSON} и \texttt{GeoPackage}).
Указанное ограничение, очевидно, приводит к излишним затратам времени на сериализацию и десериализацию данных, а так же приводит к излишнему объему требуемой памяти для хранения результатов промежуточных расчетов.

Указанный выбор в пользу файлового хранения информации был сделан в связи с необходимостью выполнения большого объема работ по тестированию результатов расчетов и необходимости их визуализации.
Локальность участков расчета выполненного в примере не позволяет указанным недостаткам перейти черту критической несовместимости с выполнением.
В следующих ревизиях программного кода упомянутые ограничения будут решены за счет использования баз данных совместимых с геопространственными данными.
